\subsection{What the finite Gaussian translation quotient does not determine}

There is a concrete finite obstruction to concluding the harmonic
square test solely from the translation group, its symplectic
action and the three orbifold inertia classes. This obstruction
is at the level of a finite peripheral model; it is not a GT
counterexample.

\begin{proposition}\label{prop:gaussian-finite-peripheral-scope}
For every odd prime \(\ell\), there is a finite quotient of the
three-punctured sphere group with peripheral orders \(4,2,4\),
and an automorphism that fixes its maximal pro-\(p\) quotient
for every prime \(p\), preserves each positive peripheral
conjugacy class, acts symplectically on its translation subgroup,
and has nontrivial Gaussian square sign.
\end{proposition}
\begin{proof}
Use the finite etale algebra
\[
 A=\mathbf F_\ell[i]=\mathbf F_\ell[t]/(t^2+1),\qquad
 J=A_{\rm add}\rtimes\langle i\rangle.
\]
The algebra has dimension two, including in the split case.
Its involution sends \(i\) to \(-i\). Write
\((v,r)(w,s)=(v+rw,rs)\), and choose
\[
 x=(0,i),\qquad y=(1,-1),\qquad z=(xy)^{-1}=(1,i).
\]
These elements have orders \(4,2,4\), respectively, and \(xyz=1\).
Their commutator is \([x,y]=(i-1,1)\).
Since \(i-1\) is invertible and conjugation by \(x\) multiplies
translations by \(i\), they generate \(J\).
Thus they define the specified quotient.

The norm-one subgroup
\[
 A^1=\{a\in A^\times:a\overline a=1\}
\]
is cyclic of order \(\ell-\chi_{-4}(\ell)\), which is divisible
by four. Choose a nonsquare \(a\in A^1\), and define
\(\phi_a(v,r)=(av,r)\). This is an automorphism.
Its determinant on the two-dimensional translation space is
\(a\overline a=1\), hence it preserves its alternating form.
It fixes \(x\). If
\[
 f=((1-a)/2,1),
\]
then \(\phi_a(y)=f^{-1}yf\). The two elements \(z=(1,i)\)
and \(\phi_a(z)=(a,i)\) are conjugate by a translation, because
translation conjugacy changes the first coordinate by
\((1-i)v\), and \(1-i\) is invertible.
Thus all three positive peripheral classes are preserved.

The maximal pro-two quotient of \(J\) is \(C_4\), and
\(\phi_a\) fixes it. Its maximal pro-\(\ell\) quotient is
trivial: killing \(i\), whose order is four, also kills every
translation by the relation \(iv=v\), since \(i-1\) is invertible.
All other single-prime quotients are trivial. Nevertheless
\[
             a^{(\ell-\chi_{-4}(\ell))/2}=-1.
\]
The affine Alexander calculation identifies its beta scalar
with \(a\): since \([x,y]=i-1\), the conjugator has coefficient
\(R=(1-a)/(2(i-1))\), and
\[
             1+(i-1)(-2)R=a.
\]
\end{proof}

This construction specifies only an automorphism of the displayed
finite quotient. It does not lift that automorphism to full GT,
impose the pentagon, or prescribe the simultaneous rational
Kummer coordinates of a lift. In particular it does not assert
that an element of \(P_D\) has nontrivial harmonic residue.
It shows why triviality of all single-prime word actions does
not by itself settle a finite quotient with translation and
rotation of different prime orders.

For a geometric attempt using the full \(\ell\)-division
translation group on a torus, the multiplication-by-\(\ell\)
cover is characteristic on closed torus homology. Its liftable
mapping-class subgroup is therefore the full mapping-class
group. This removes the particular normalizer-surjectivity
obstruction of a cyclic isogeny. On first homology, however,
the resulting intertwining is initially just commutation with
the scalar \(\ell I\). The finite proposition shows that the
translation and orbifold inertia data still permit both
Gaussian square classes. A proof of the desired relation
must additionally identify the tangential or rotation
correction with the divided Kummer coordinate
\(\rho_\ell^I/2\bmod2\). That identification remains open here.
